\documentclass[11pt]{article}

\usepackage[margin=1in]{geometry}
\usepackage{listings}
\usepackage{xcolor}
\usepackage{hyperref}
\usepackage{amsmath}
\usepackage{parskip}

\hypersetup{
  colorlinks=true,
  linkcolor=blue!60!black,
  urlcolor=blue!60!black
}

\lstset{
  basicstyle=\ttfamily\small,
  frame=single,
  breaklines=true,
  showstringspaces=false,
  columns=fullflexible,
  keywordstyle=\color{blue!70!black},
  commentstyle=\color{green!40!black},
  stringstyle=\color{red!60!black}
}

\title{CS201 Winter 2025 --- Lecture Notes}
\author{Clarissa Littler}
\date{Winter 2025}

\begin{document}
\maketitle

\section*{Preface}

These are my running notes for the term --- one section per class
meeting, walking through whatever we covered that day with the
actual code we typed up in class. The source files all live in the
\texttt{lectureN/} subdirectories of the repo, and the snippets I
quote here are pulled straight out of those files. If you missed a
lecture, want a refresher, or are studying for the final, this is
the one-stop overview.

\tableofcontents

\section{Lecture 1 --- Welcome, What an OS Is, and Bits}

This was the orientation lecture. I sketched out what we're going to
spend the term on, made a pitch for why systems programming is
interesting, and then we dove straight into the bit-level stuff that
underlies everything else.

\subsection*{What this class is}

The class is a mix of:
\begin{itemize}
  \item How operating systems actually work
  \item How programs run on a real machine
  \item C programming
  \item x86-64 assembly
  \item The contingencies of doing all of this on Linux
  \item Concurrent and parallel programming
\end{itemize}
Plus some smaller tangents along the way --- a tiny bootloader, some
CTF-flavored challenges, me rambling about microkernels, etc.

\subsection*{What is an operating system, anyway?}

We took a few minutes to brainstorm. The short version: an OS is the
piece of software that mediates between hardware and everything else,
and the things you tend to want from it are process management,
memory management, file I/O, and a uniform interface to devices.

\subsection*{A quick C tour}

Most of you have C++ experience, so I focused on what's different in
C. The minimum-viable program is just:

\begin{lstlisting}[language=C]
#include <stdio.h>

int main(){
  printf("Hello, world!\n");
  return 0;
}
\end{lstlisting}

A few things to call out coming from C++:
\begin{itemize}
  \item I/O is \texttt{printf}/\texttt{scanf}, not streams. The
    format specifiers (\texttt{\%d}, \texttt{\%s}, \texttt{\%lu},
    \ldots) are how the function knows what types to expect.
  \item No pass-by-reference --- you pass pointers if you want a
    callee to modify a caller's variable.
  \item No \texttt{cin}-style line-buffered input convenience;
    \texttt{scanf} reads the format you tell it to.
  \item \texttt{malloc}/\texttt{free} instead of \texttt{new}/\texttt{delete}.
\end{itemize}

\texttt{adder.c} drove these home: read two integers and print their
sum.

\begin{lstlisting}[language=C]
int num1; int num2;
printf("Enter a number: ");
scanf("%d",&num1);
printf("Enter another number: ");
scanf("%d",&num2);
printf("So %d and %d added together make %d\n",num1,num2,num1+num2);
\end{lstlisting}

The \texttt{\&num1} is the part that matters: \texttt{scanf} needs an
\emph{address} to store into, because there's no other way to write
out-parameters in C.

\texttt{echo.c} is similar but for a string; it also shows that an
array name is essentially a pointer:

\begin{lstlisting}[language=C]
char msg[50];
printf("Type somethin' will ya: ");
scanf("%s",msg);
printf("You said: %s\n",msg);
printf("And msg is actually: %p\n",msg);
\end{lstlisting}

\subsection*{\texttt{malloc} and structs}

Two tiny demos. \texttt{malloc1.c} is just enough to show that
\texttt{malloc} takes a byte count and returns a pointer:

\begin{lstlisting}[language=C]
int* arr = malloc(10000000*sizeof(int));
free(arr);
\end{lstlisting}

\texttt{malloc2.c} adds a struct and shows the two field-access
syntaxes:

\begin{lstlisting}[language=C]
struct Point* pointy = malloc(sizeof(struct Point));
(*pointy).x = 10;
(*pointy).y = 20;
pointy->x = 50;     // shorthand for (*pointy).x
pointy->y = 40;
free(pointy);
pointy = 0;
\end{lstlisting}

The \texttt{->} arrow is just sugar for ``dereference and access a
field.'' Always free what you malloc, and zero out the pointer
afterward so you don't accidentally use it again.

\subsection*{Bits and bytes}

A byte is 8 bits. The bits in a byte read most-significant on the
left, so \texttt{10010101} unpacked is
\[
  1 \cdot 2^7 + 0 \cdot 2^6 + 0 \cdot 2^5 + 1 \cdot 2^4
  + 0 \cdot 2^3 + 1 \cdot 2^2 + 0 \cdot 2^1 + 1 \cdot 2^0 = 149.
\]

The bitwise operators in C are:
\begin{itemize}
  \item \texttt{\&} --- bitwise AND.
  \item \texttt{|} --- bitwise OR.
  \item \texttt{\^{}} --- bitwise XOR.
  \item \texttt{\textasciitilde{}} --- bitwise NOT.
\end{itemize}
Don't confuse them with the boolean operators \texttt{\&\&} and
\texttt{||}, which short-circuit and produce 0/1.

\texttt{bit1.c} is the first ``play with bits'' program: read an
integer, print all 32 bits of it, ask which bit to flip, flip it,
loop forever:

\begin{lstlisting}[language=C]
void printBits(int n){
  for(int i = 31; i >= 0; i--){
    printf("%d",(n >> i)&1);
  }
  printf("\n");
}

void flipBit(int* n, int c){
  *n = (*n) ^ (1 << c);
}
\end{lstlisting}

The XOR trick is the key thing: XOR-ing with 1 flips a bit, XOR-ing
with 0 leaves it alone. So \texttt{n \^{} (1 << c)} flips exactly the
$c$th bit.

\texttt{bitshiftDemo.c} is the most boring file in the directory but
its one line ---
\begin{lstlisting}[language=C]
printf("%d\n",1024 >> 2);
\end{lstlisting}
--- captures something I want you to internalize: shifting right by
$k$ divides by $2^k$ (for non-negative integers), and shifting left
by $k$ multiplies by $2^k$.

\subsection*{Two's complement}

If we have $n$ bits we can represent unsigned $0$ through $2^n - 1$
with the usual binary expansion. To represent signed integers,
we flip the sign of the top bit's weight:
\[
  -b_{n-1} \cdot 2^{n-1} + \sum_{i=0}^{n-2} b_i \cdot 2^i.
\]
So \texttt{11111111} as a signed 8-bit number is
$-128 + 64 + 32 + 16 + 8 + 4 + 2 + 1 = -1$. Negating a number is
``flip all the bits and add 1,'' or equivalently ``subtract 1 then
flip all the bits.''

We'll come back to this next time and play with the actual
representations of negative numbers and floats.

\section{Lecture 2 --- Sizes, Pointers, Endianness, and IEEE-754}

Today we made data types concrete: how big they are, what their
addresses look like, how a multi-byte value is stored on this
machine, and finally how floats work.

\subsection*{Sizes}

\texttt{size1.c} is just \texttt{sizeof} on the basic types:

\begin{lstlisting}[language=C]
printf("The size of an int is: %lu\n",sizeof(int));
printf("The size of a char is: %lu\n",sizeof(char));
printf("The size of a float is: %lu\n",sizeof(float));
printf("The size of a double is: %lu\n",sizeof(double));
printf("The size of a long int is: %lu\n",sizeof(long int unsigned));
\end{lstlisting}

On our 64-bit Linux boxes: int = 4, char = 1, float = 4, double = 8,
long = 8.

\subsection*{Is char signed?}

\texttt{charSigned.c} settles a perennial question:

\begin{lstlisting}[language=C]
char test = -1;
if(test > 0){
  printf("Well, char is unsigned!\n");
}
else{
  printf("char is signed, I guess\n");
}
\end{lstlisting}

The C standard does not pin down whether plain \texttt{char} is
signed or unsigned --- it's \emph{implementation-defined}. On x86-64
GCC it's signed, so this prints ``char is signed, I guess.'' If you
want to be sure, write \texttt{signed char} or \texttt{unsigned
char}.

\subsection*{Pointer arithmetic}

\texttt{pointArith.c} prints the addresses of consecutive elements of
two arrays of different types:

\begin{lstlisting}[language=C]
int arr1[5] = {0,1,2,3,4};
double arr2[5] = {0,1,2,3,4};
for(int i = 1; i < 5; i++){
  printf("The %dth element of arr1 is at: %p\n",i,(void*)(arr1 + i));
  printf("The %dth element of arr2 is at: %p\n",i,(void*)(arr2 + i));
}
\end{lstlisting}

The takeaway is that pointer arithmetic scales with the size of the
\emph{pointee}. Adding 1 to an \texttt{int*} advances the address by
4 bytes; adding 1 to a \texttt{double*} advances by 8. The compiler
quietly multiplies your index by \texttt{sizeof(*p)} on every
\texttt{p + i}.

\subsection*{Endianness}

\texttt{whichEnd.c} pokes at how a multi-byte integer is actually
laid out in memory:

\begin{lstlisting}[language=C]
int num = 0x12345678;
unsigned char* pointy = (unsigned char*)&num;
printf("Our number's bytes, in order, are: ");
for(size_t i = 0; i < sizeof(num); i++){
  printf(" %x ",*(pointy+i));
}
\end{lstlisting}

The cast to \texttt{unsigned char*} is the trick: now pointer
arithmetic moves us one byte at a time, so we can walk the actual
storage. On x86-64 you'll see the bytes come out as \texttt{78 56 34
12} --- reversed from how we wrote the literal. That's
little-endian: least-significant byte first. Big-endian machines
would print it in the order it was written. Network byte order is
big-endian, which is why \texttt{htons} and friends exist.

\subsection*{IEEE-754}

I wrote out the 32-bit float layout on the board:
\begin{itemize}
  \item 1 bit sign
  \item 8 bits exponent (biased by 127)
  \item 23 bits significand (the ``mantissa'')
\end{itemize}

For normalized numbers the value is
\[
  (-1)^s \cdot 2^{|e| - 127} \cdot (1 + |m|),
\]
where the significand is read as a sum of \emph{negative} powers of
2: bit position $k$ after the implicit point contributes $2^{-k}$.

A worked example: \texttt{0 10000001 1100\ldots0} is
\[
  (-1)^0 \cdot 2^{129 - 127} \cdot (1 + \tfrac12 + \tfrac14)
  = 4 \cdot \tfrac{7}{4} = 7.
\]
The reason there's an implicit leading 1 is that the significand is
required to be $\ge 1$ and $< 2$ in the normalized case, so we save
a bit by not storing it.

The edge cases are handled by the exponent:
\begin{itemize}
  \item Exponent all 0s, significand 0 $\Rightarrow$ $\pm 0$.
  \item Exponent all 0s, significand $\ne 0$ $\Rightarrow$
    subnormal: $(-1)^s \cdot 2^{-126} \cdot |m|$ (no implicit 1).
  \item Exponent all 1s, significand 0 $\Rightarrow$ $\pm \infty$.
  \item Exponent all 1s, significand $\ne 0$ $\Rightarrow$
    \texttt{NaN}.
\end{itemize}

\texttt{floatMunger.c} is just \texttt{bit1.c} from last time, but
the integer-bit-flipper is run on a \texttt{float} via a pointer
cast:

\begin{lstlisting}[language=C]
float num;
scanf("%f",&num);
while(1){
  printf("The number is: %.40f\n",num);
  printBits((int*)&num);
  printf("Which bit to flip?: ");
  scanf("%d",&choice);
  flipBit((int*)&num,choice);
}
\end{lstlisting}

This is the single best way to internalize the format: type a
familiar number, see its bits, flip the sign bit, watch it negate;
flip the low exponent bit, watch it double.

\subsection*{Where floats lie to you}

\texttt{floatTest.c} shows the canonical disappointment:

\begin{lstlisting}[language=C]
double num1 = 0;
double num2 = 0.001;
for(int i = 0; i < 10000; i++){
  num1 = num1 + num2;
}
if(num1 == 10){
  printf("Yee-caw!\n");
}
else {
  printf("Whoops!\n");
  printf("num1 is really: %.20f\n",num1);
}
\end{lstlisting}

\texttt{0.001} is not exactly representable in binary floating point.
The error is small per addition, but we accumulate it 10{,}000 times,
and the final answer is not 10. This is the kind of thing where
``check floats with \texttt{==}'' will burn you forever.

The deeper reason it's unavoidable: there are uncountably many real
numbers but only finitely many floats, so we're forced to round most
reals to the nearest representable one. IEEE-754 specifies
``round-to-even'' on ties, so the rounding doesn't bias one direction
on average.

\subsection*{Posits, briefly}

I previewed posits because they're the topic of the next discussion
and they're the next-day reading. The short version: they're an
alternative number format with variable-width fields, parameterized
by total bit width $n$ and max exponent bits $es$. They give up some
precision at the extremes in exchange for more precision in the
middle of the number line, and they have just one ``not-a-real''
sentinel instead of IEEE's menagerie. The tutorial in
\texttt{positTutorial.org} walks through the encoding in detail.

\section{Lecture 3 --- Review, Stacks, and First Assembly}

Today was the pivot from C to assembly. We reviewed IEEE and posits,
talked about the upcoming stack assignment, and then dove into the
GNU assembler.

\subsection*{IEEE-754 review}

I worked one more example on the board to nail down the formula. For
the normalized 32-bit case,
\[
  (-1)^{|s|} \cdot 2^{|e| - 127} \cdot (1 + |m|).
\]
Take \texttt{1 00111111 10100\ldots0}: sign 1, exponent $63$ so
$|e| - 127 = -64$, significand $1/2 + 1/8 = 5/8$, so the value is
$-1 \cdot 2^{-64} \cdot (13/8)$, a tiny negative number.

I also pointed out a subtlety: if we allowed the significand to be
$\ge 2$, a single value would have multiple representations
(shift the significand right and bump the exponent). Forcing the
implicit leading 1 makes the encoding unique.

\subsection*{Posits, more detail}

Same idea as last time but written out more carefully. With $n = 8$,
$es = 2$, the bits split as: 1 sign bit, then \emph{regime} bits
counted by run-length, then up to $es$ exponent bits, then whatever's
left for the significand. The value is
\[
  (-1)^s \cdot 2^{2^{es} \cdot k} \cdot 2^{|e|} \cdot (1 + |m|),
\]
where $k$ is determined by the regime: if the regime starts with 1
and has length $r$, then $k = r - 1$; if it starts with 0, then
$k = -r$. Worked example: \texttt{01101101} parses as
\[
  1 \cdot 2^{2^2} \cdot 2^3 \cdot (1 + 1/4)
  = 16 \cdot 8 \cdot 5/4 = 160.
\]

If you want to see the full picture --- including the
``not-a-real'' sentinel and how to do arithmetic --- read
\texttt{positTutorial.org} in the repo.

\subsection*{The stack assignment}

Reminder: a stack is a LIFO data structure with two operations,
\texttt{push} and \texttt{pop}. For the assignment I want you to
implement one in C: pick either a linked list or an array-with-length
representation, expose
\begin{lstlisting}[language=C]
void push(struct Stack* s, int x);
int pop(struct Stack* s);
\end{lstlisting}
and write a tiny driver that demonstrates LIFO order.

\subsection*{Compiling to assembly: \texttt{gcc -S}}

The first thing we did with assembly was to look at what GCC
generates for our \texttt{hello.c}. \texttt{gcc -S hello.c} dumps
\texttt{hello.s}:

\begin{lstlisting}
        .file	"hello.c"
        .intel_syntax noprefix
        .text
        .section	.rodata
.LC0:
        .string	"Hello CS201!"
        .text
        .globl	main
        .type	main, @function
main:
        endbr64
        push    rbp
        mov     rbp, rsp
        lea     rax, .LC0[rip]
        mov     rdi, rax
        call    puts@PLT
        mov     eax, 0
        pop     rbp
        ret
\end{lstlisting}

A few things to notice:
\begin{itemize}
  \item There are \emph{sections}: \texttt{.text} for code,
    \texttt{.rodata} for read-only data (the string lives there).
  \item GCC swapped our \texttt{printf} for \texttt{puts}, because
    the format string had no format specifiers --- a tiny optimization.
  \item This file is in Intel syntax (\texttt{.intel\_syntax
    noprefix}); the GNU assembler defaults to \emph{AT\&T} syntax,
    which is what we'll use the rest of the term. PSU does AT\&T,
    PCC does AT\&T, that's the convention.
\end{itemize}

The big visible difference between the two syntaxes is operand order:
in AT\&T, \texttt{mov S, D} means ``move $S$ to $D$'' (source first);
in Intel, \texttt{mov D, S} means the same thing (destination first).
Register names get a \texttt{\%} prefix in AT\&T, immediates get a
\texttt{\$} prefix.

\subsection*{Our first hand-written assembly: \texttt{ret1.s}}

\begin{lstlisting}
        .text
        .global _start

_start:
        mov $60,%rax
        mov $-1,%rdi
        syscall
\end{lstlisting}

This is the smallest assembly program that does something
observable. \texttt{\_start} is the entry point that the linker
expects (no C runtime to call \texttt{main} for us). We put 60 in
\texttt{\%rax} (the syscall number for \texttt{exit} on x86-64
Linux), put our exit status in \texttt{\%rdi}, and execute the
\texttt{syscall} instruction. The shell will report the exit status
as 255, because exit codes are masked to 8 bits.

To assemble and link this by hand:
\begin{lstlisting}
as -o ret1.o ret1.s
ld -o ret1 ret1.o
\end{lstlisting}
Run it, then \texttt{echo \$?}. That's our entire workflow for the
next few lectures.

\section{Lecture 4 --- Instructions, Data, Compares, and Loops}

Today we filled out the assembly toolkit: arithmetic instructions
with their size suffixes, the \texttt{.data} section for globals,
the addressing modes for arrays, and \texttt{cmp}/\texttt{jmp} for
control flow.

\subsection*{Instruction sizes}

Most arithmetic instructions come in size-suffixed flavors:
\begin{itemize}
  \item \texttt{q} --- quad, 64 bits
  \item \texttt{l} --- long, 32 bits
  \item \texttt{w} --- word, 16 bits
  \item \texttt{b} --- byte, 8 bits
\end{itemize}
You can usually let the assembler infer the size from the operands,
but sometimes you need to be explicit (e.g.\ when both operands are
immediates or when you're operating on memory through a pointer with
no register involved).

\subsection*{\texttt{add} and \texttt{sub}}

\texttt{add2.s}:

\begin{lstlisting}
_start:
        mov $1,%rax
        mov $2,%rdx
        sub %rax,%rdx           # sub s,d -> d = d - s
        mov %rdx,%rdi
        mov $60,%rax
        syscall
\end{lstlisting}

In AT\&T syntax, \texttt{add S, D} computes \texttt{D = D + S}, and
\texttt{sub S, D} computes \texttt{D = D - S}. The destination is on
the right. This will trip you up at some point this term --- I
guarantee it. The result of \texttt{2 - 1} is 1, which we move into
\texttt{\%rdi} so it becomes the exit status.

\subsection*{\texttt{cmp} and \texttt{jmp}}

\texttt{cmp S, D} sets the CPU flags based on \texttt{D - S} but
discards the result. The conditional jumps then branch on those
flags. \texttt{cmp1.s}:

\begin{lstlisting}
_start:
        mov $1,%rax
        mov $2,%rbx

        cmp %rbx,%rax           # compare as if %rax - %rbx
        jne less
        mov $-1,%rdi
        jmp finish
less:
        mov $10,%rdi
finish:
        mov $60,%rax
        syscall
\end{lstlisting}

The cheat sheet for the conditional jumps:
\begin{itemize}
  \item \texttt{je} / \texttt{jne} --- equal / not equal.
  \item \texttt{jl} / \texttt{jle} --- less / less-or-equal (signed).
  \item \texttt{jg} / \texttt{jge} --- greater / greater-or-equal
    (signed).
  \item \texttt{jb} / \texttt{ja} --- below / above (unsigned).
\end{itemize}
Pay attention to the operand-order quirk: \texttt{cmp \%rbx,\%rax}
followed by \texttt{jl} branches if $\%rax < \%rbx$ in the AT\&T
mental model. (I read it as ``D fewer than S'' to keep myself
straight.)

\subsection*{The \texttt{.data} section}

To declare globals you put them in \texttt{.data}, with directives
that pick the size:
\begin{itemize}
  \item \texttt{.byte} --- 8 bits
  \item \texttt{.word} --- 16 bits
  \item \texttt{.long} --- 32 bits
  \item \texttt{.quad} --- 64 bits
\end{itemize}
You can list multiple values: \texttt{.long 1,2,3,4}.

\texttt{data1.s} loads a single global, increments it, and exits
with its value:

\begin{lstlisting}
        .section .data
num1:   .long 200

        .section .text
        .global _start

_start:
        mov num1(%rip),%rbx
        addl $10,(%rbx)
        mov num1,%rdi
        mov $60,%rax
        syscall
\end{lstlisting}

\textbf{Heads up:} this file as written has a couple of subtleties
worth flagging. \texttt{num1} is a \texttt{.long} (32 bits), but the
first \texttt{mov num1(\%rip),\%rbx} loads 64 bits into a 64-bit
register --- the upper half is whatever happens to live next to
\texttt{num1} in the \texttt{.data} section, which is fine here only
because the program exits immediately. The cleaner version would use
\texttt{movl num1(\%rip),\%ebx} to load the 32-bit value into the
32-bit register \texttt{\%ebx} (which auto-zeros the upper half).
Mixing \texttt{.long} with the 64-bit register names is the single
most common bug pattern I see all term --- watch for it in your own
code.

\subsection*{Arrays and addressing modes}

\texttt{data2.s} is a for-loop that sums an array:

\begin{lstlisting}
        .section .data
arr1:   .long 1,2,3,4

        .section .text
        .global _start

_start:
        lea arr1(%rip),%rbx     # %rbx = &arr1
        mov $0,%rcx             # i = 0
        mov $0,%rax             # sum = 0

loopBody:
        cmp $4,%rcx
        jge done
        add (%rbx,%rcx,4),%rax  # sum += arr1[i]
        add $1,%rcx
        jmp loopBody

done:
        mov %rax,%rdi
        mov $60,%rax
        syscall
\end{lstlisting}

The new piece is the addressing mode \texttt{(\%rbx,\%rcx,4)}, which
computes the address $\texttt{\%rbx} + \texttt{\%rcx} \cdot 4$. The
syntax is \texttt{(base, index, scale)} with scale in \{1, 2, 4, 8\}.
This is exactly how a C-style array index compiles, and you'll see
it constantly.

\textbf{Same gotcha as before:} \texttt{arr1} is \texttt{.long}
(4-byte) elements, so the scale is 4 --- but \texttt{add
(\%rbx,\%rcx,4),\%rax} does a 64-bit load into \texttt{\%rax}, which
reads 8 bytes per iteration. The two adjacent \texttt{.long} values
fit into one 8-byte read, so on the first iteration we add elements
0 and 1 together, on the next iteration elements 1 and 2, etc. The
``correct'' version would use \texttt{addl (\%rbx,\%rcx,4),\%eax} ---
32-bit operand size --- so each load picks up exactly one
array element. Try both and watch the exit status change.

\texttt{data3.s} is the same loop, but it factors out the address
computation using \texttt{lea}:

\begin{lstlisting}
        lea (%rbx,%rcx,4),%rdx  # %rdx = &arr1[i]; no dereference
        add (%rdx),%rax         # sum += *(%rdx)
\end{lstlisting}

The mnemonic \texttt{lea} is for ``load effective address'': it
computes the address that the addressing mode would point to and
stuffs that into the destination, but does \emph{not} dereference.
Useful when you want a pointer to something rather than the thing
itself --- the assembly equivalent of C's \texttt{\&}.

\subsection*{Hello world preview}

I sketched the structure of the next program: \texttt{.asciz} for a
null-terminated string in \texttt{.data}, the \texttt{message\_len =
. - message} idiom for getting the length at assembly time, and the
\texttt{write} syscall (number 1, args fd / buf / count in
\texttt{\%rdi} / \texttt{\%rsi} / \texttt{\%rdx}). We'll write the
full program first thing next class.

\section{Lecture 5 --- Hello, Echo, Labels, and \texttt{call}/\texttt{ret}}

Today we wrote ``hello, world'' end-to-end, then \texttt{echo}
end-to-end, then introduced labels and the \texttt{call}/\texttt{ret}
mechanism that lets us write actual functions.

\subsection*{Reference: the calling convention}

Quick reminder of the registers and how arguments get passed:
\begin{itemize}
  \item Argument registers (in order): \texttt{\%rdi}, \texttt{\%rsi},
    \texttt{\%rdx}, \texttt{\%rcx}, \texttt{\%r8}, \texttt{\%r9}.
  \item Return value: \texttt{\%rax}.
  \item Caller-saved: \texttt{\%rax}, \texttt{\%rcx}, \texttt{\%rdx},
    \texttt{\%rsi}, \texttt{\%rdi}, \texttt{\%r8}--\texttt{\%r11}.
  \item Callee-saved: \texttt{\%rbx}, \texttt{\%rbp},
    \texttt{\%r12}--\texttt{\%r15}.
\end{itemize}
Syscall arguments use a slightly different layout (4th arg is
\texttt{\%r10} not \texttt{\%rcx}), but for our purposes the syscalls
we use only need three or fewer args, so we don't have to think
about it.

Syscall numbers we'll be using:
\begin{itemize}
  \item \texttt{0} --- \texttt{read(fd, buf, count)}
  \item \texttt{1} --- \texttt{write(fd, buf, count)}
  \item \texttt{60} --- \texttt{exit(status)}
\end{itemize}
And the file descriptors: \texttt{0} = stdin, \texttt{1} = stdout,
\texttt{2} = stderr.

\subsection*{Hello, world}

\texttt{hello.s}:

\begin{lstlisting}
        .section .rodata
msg:    .asciz "Hello world!\n"
hellolen = . - msg

        .section .text
        .globl _start

_start:
        mov $1,%rax             # syscall: write
        mov $2,%rdi             # fd = 2 (stderr)
        lea msg(%rip),%rsi
        mov $hellolen,%rdx
        syscall

        xor %rdi,%rdi
        mov $60,%rax
        syscall
\end{lstlisting}

A few points:
\begin{itemize}
  \item \texttt{.asciz} lays down the characters followed by a null
    terminator.
  \item \texttt{hellolen = . - msg} is the canonical idiom for
    string length: \texttt{.} is the current address, so subtracting
    \texttt{msg} gives the byte length (including the newline and
    the trailing \texttt{\textbackslash 0}).
  \item I deliberately wrote \texttt{fd = 2} (stderr) here to make
    the point that fd 0 / 1 / 2 are conventions, not magic; you can
    write to any of them. Change it to \texttt{\$1} for stdout.
  \item \texttt{xor \%rdi,\%rdi} is the idiomatic way to zero a
    register --- shorter encoding than \texttt{mov \$0, \%rdi}.
\end{itemize}

\subsection*{Echo, with the \texttt{.bss} section}

To read input we need somewhere to read \emph{into}. The
\texttt{.bss} section is for uninitialized, zero-filled storage
that costs nothing on disk:

\begin{lstlisting}
        .section .bss
buff:   .skip 128
\end{lstlisting}

\texttt{.skip 128} reserves 128 bytes at the label \texttt{buff}.

\texttt{echo.s} reads from stdin and writes back what it got:

\begin{lstlisting}
        .section .bss
buff:   .skip 128

        .section .text
        .globl _start

_start:
        mov $0,%rax             # read
        mov $1,%rdi             # stdin
        lea buff(%rip),%rsi
        mov $128,%rdx
        syscall

        mov %rax,%rdx           # bytes-read becomes byte count
        mov $1,%rax             # write
        lea buff(%rip),%rsi
        mov $0,%rdi             # stdout
        syscall

        xor %rdi,%rdi
        mov $60,%rax
        syscall
\end{lstlisting}

The clever bit is \texttt{mov \%rax,\%rdx} between the two syscalls:
\texttt{read} returns the number of bytes actually read in
\texttt{\%rax}, and we feed that straight into \texttt{write}'s
byte-count argument. Otherwise we'd echo the uninitialized tail of
the buffer too.

\textbf{Bug worth flagging:} the file as written has the
\texttt{\%rdi} values for read/write swapped --- it sets
\texttt{\%rdi = 1} for the \texttt{read} (which would mean ``read
from stdout''!) and \texttt{\%rdi = 0} for the \texttt{write} (which
would mean ``write to stdin''). On a normal interactive terminal
that fd 0 is opened read/write, so it ``works'' (it reads from your
terminal and writes back to it), but the right values are 0 for
read and 1 for write. This is the same fd-mixup pattern I'll keep
flagging all term.

\subsection*{Conditional echoes: \texttt{echoBroken.s}}

\texttt{echoBroken.s} branches based on what's in the buffer:

\begin{lstlisting}
        .section .bss
buff:   .skip 128

        .section .rodata
msg1:   .asciz "Was zero\n"
msg1len = . - msg1
msg2:   .asciz "Not zero\n"
msg2len = . - msg2

        .section .text
        .globl _start

_start:
        lea buff(%rip),%rbx
        cmpb $0,(%rbx)
        je isZero
        jmp notZero

isZero:
        lea msg1(%rip),%rsi
        mov $msg1len,%rdx
        jmp writeLabel

notZero:
        lea msg2(%rip),%rsi
        mov $msg2len,%rdx

writeLabel:
        mov $1,%rax
        mov $0,%rdi
        syscall
        xor %rdi,%rdi
        mov $60,%rax
        syscall
\end{lstlisting}

The point isn't really the conditional --- the ``broken'' part is
that we never actually \emph{read} into \texttt{buff}, so it's
always zero. (The corresponding \texttt{read} block is commented out
at the top.) Same fd-1/fd-0 swap on the write call, too.

The new instruction here is \texttt{cmpb \$0,(\%rbx)} --- compare the
single byte at address \texttt{\%rbx} against \texttt{0}. The
\texttt{b} suffix is required because there's no register to hint at
the size; we need to tell the assembler we mean a byte.

\subsection*{Local labels and \texttt{call}/\texttt{ret}}

\texttt{labelTest.s} introduces two new ideas: numeric local labels
(\texttt{1:}, \texttt{2:}) and the \texttt{call}/\texttt{ret} pair.

\begin{lstlisting}
_start:
        call biglabel2
        call biglabel1

1:
        mov $60,%rax
        xor %rdi,%rdi
        syscall

biglabel1:
        jmp 1f
1:
        lea msg1(%rip),%rsi
        mov $msg1len,%rdx
        mov $1,%rax
        mov $0,%rdi
        syscall
        ret

biglabel2:
        jmp 1f
1:
        lea msg2(%rip),%rsi
        mov $msg2len,%rdx
        mov $1,%rax
        mov $0,%rdi
        syscall
        ret
\end{lstlisting}

Two new mechanics:

\begin{itemize}
  \item \emph{Numeric local labels.} \texttt{1:} can be reused as
    many times as you want. \texttt{jmp 1f} means ``jump to the next
    \texttt{1:} forward,'' \texttt{jmp 1b} means ``jump to the
    nearest \texttt{1:} backward.'' Saves you inventing names for
    one-off jumps inside a function.
  \item \texttt{call}/\texttt{ret}. \texttt{call label} pushes the
    address of the next instruction and jumps to \texttt{label};
    \texttt{ret} pops that address and jumps back. So now we have
    real subroutines that can be called from multiple places.
\end{itemize}

The fd-0 vs fd-1 mix-up is in this file too. (When I notice these
later when teaching from the file, I treat them as ``what does this
do? what should it do?'' moments --- they're more useful as
debugging exercises than mistakes to fix silently.)

\section{Lecture 6 --- Functions with Stack Frames}

Today we made functions feel like real functions: they take
arguments, they have local variables on the stack, and they return
values cleanly. The two files in this lecture both compute something
silly so we can focus on the mechanics.

\subsection*{Plan}

The arc:
\begin{enumerate}
  \item Push the old base pointer.
  \item Move \texttt{\%rsp} into \texttt{\%rbp} so the new frame
    starts where the stack currently is.
  \item Subtract some multiple of 16 from \texttt{\%rsp} to carve
    out room for locals.
  \item Use \texttt{-8(\%rbp)}, \texttt{-16(\%rbp)}, etc.\ as your
    local variables.
  \item On the way out, restore \texttt{\%rsp}, pop \texttt{\%rbp},
    and \texttt{ret}.
\end{enumerate}

Why a multiple of 16? Because the ABI requires \texttt{\%rsp} to be
16-byte aligned at the point of every \texttt{call}, and some
instructions (notably the SSE ones that shuffle 128-bit values) will
fault if it isn't. Better to over-allocate by 8 bytes than to debug
a mysterious crash inside a library function.

\subsection*{\texttt{fun1.s}: a function without locals}

\begin{lstlisting}
        ## int fun(int a, int b){
        ##   return 3*(a + b);
        ## }
fun:
        mov $0,%rax
        add %rdi,%rax
        add %rsi,%rax
        imul $3,%rax            # signed multiply
        ret

_start:
        mov $2,%rdi
        mov $3,%rsi
        call fun
        mov %rax,%rdi
        mov $60,%rax
        syscall
\end{lstlisting}

Two notes:
\begin{itemize}
  \item Arguments come in via \texttt{\%rdi}, \texttt{\%rsi}, the
    return value goes back in \texttt{\%rax}.
  \item \texttt{imul} is signed multiply; \texttt{mul} is unsigned.
    Use \texttt{imul} unless you know your operands are unsigned.
\end{itemize}

This function has no locals, so we don't need to set up a stack
frame. Result: $3 \cdot (2 + 3) = 15$, so the program exits with
status 15.

\subsection*{\texttt{fun2.s}: a function with locals}

\begin{lstlisting}
        ## long fun(long a, long b){
        ##   long c = a + b;
        ##   long d = 2 + b;
        ##   return c + d;
        ## }
fun:
        push %rbp
        mov %rsp,%rbp
        sub $16,%rsp           # room for two quadword locals

        mov %rdi,-8(%rbp)      # local c = a
        add %rsi,-8(%rbp)      # c += b
        mov %rsi,-16(%rbp)     # local d = b
        addq $2,-16(%rbp)      # d += 2
        mov -16(%rbp),%rax
        add -8(%rbp),%rax      # %rax = d + c

        leave                  # mov %rbp,%rsp; pop %rbp
        ret

_start:
        mov $3,%rdi
        mov $4,%rsi
        call fun
        mov %rax,%rdi
        mov $60,%rax
        syscall
\end{lstlisting}

This is the canonical shape of an assembly function from now on.
Walk through it once on paper:
\begin{itemize}
  \item \texttt{push \%rbp} stashes the caller's base pointer
    on the stack.
  \item \texttt{mov \%rsp,\%rbp} says ``the new frame starts here.''
  \item \texttt{sub \$16,\%rsp} reserves 16 bytes of locals; we
    treat those as two 8-byte slots at \texttt{-8(\%rbp)} and
    \texttt{-16(\%rbp)}.
  \item \texttt{leave} is shorthand for ``\texttt{mov \%rbp,\%rsp;
    pop \%rbp},'' undoing the prologue in two cycles.
\end{itemize}

With \texttt{a = 3, b = 4}: \texttt{c = 7, d = 6}, return \texttt{c
+ d = 13}. Exit status 13.

\subsection*{Aside: closures}

I left \texttt{closure.lisp} in the directory as a tiny aside about
why the stack alone isn't enough for some languages:

\begin{lstlisting}
(defun closure-test (y)
  (let ((beep y))
    #'(lambda (x)
        (incf beep)
        (+ x beep))))
\end{lstlisting}

The lambda closes over \texttt{beep} --- it has to keep
\texttt{beep} alive past the return of \texttt{closure-test}. A
plain stack frame can't support that, because the frame is gone the
moment the function returns. Languages with closures either heap-allocate
the closed-over values or use a more elaborate environment representation.
That's the entry point to a much longer conversation about runtimes,
which we won't dive into in this class --- but
\texttt{assemblyGuide.org} in the repo has more on what your C
compiler is doing under the hood.

\section{Lecture 7 --- Forking, Signals, and a Loose End in Assembly}

We swing back to C now, with two big topics: \texttt{fork()} and
signals. There's also one piece of leftover assembly business at
the top.

\subsection*{Loose end: \texttt{leaQuestion.s}}

Someone asked about \texttt{lea} after lecture, so here's a tiny
file that uses it:

\begin{lstlisting}
        .section .data
arr:    .long 1,2,3,4

fun1:
        push %rbp
        mov %rsp,%rbp
        sub $16,%rsp
        movl $3,-16(%rbp)
        mov -16(%rbp),%rcx
1:
        cmp $0,%rcx
        jl 2f
        lea arr(%rip),%rbx
        movl $0,(%rbx,%rcx,4)
        dec %rcx
        jmp 1b
2:
        leave
        ret

_start:
        call fun1
        mov $60,%rax
        mov arr(%rip),%rdi
        syscall
\end{lstlisting}

This zeros out the four \texttt{.long} elements of \texttt{arr}
counting down from 3, then exits with whatever \texttt{arr[0]} is
(zero, post-clear). The point of the file is the
\texttt{(\%rbx,\%rcx,4)} addressing: base + index * 4, just like in
C indexing.

\subsection*{Process memory layout}

Before forks, a tour of how Linux lays out a process in memory ---
text segment at the bottom, then \texttt{.rodata} / \texttt{.data} /
\texttt{.bss}, then the heap growing upward, the stack growing
downward from a high address, with mmap'd regions in between. The
diagrams at \url{https://simonis.github.io/Memory/} are useful if
you want to see the actual addresses.

\subsection*{\texttt{fork()}}

\texttt{fork()} clones the current process. After it returns, you
have two processes, both at the next instruction, both with copies of
all your variables.

\texttt{pid1.c} just shows you who you are:

\begin{lstlisting}[language=C]
printf("Hello! My id is %d and my parent's id is %d\n",
       getpid(), getppid());
\end{lstlisting}

\texttt{pid2.c} actually forks:

\begin{lstlisting}[language=C]
pid_t pid = fork();
if (pid < 0){
  perror("Fork failed");
  return 1;
}
else if (pid == 0){
  printf("Hello from the child process! My PID is %d\n", getpid());
}
else {
  printf("Hello from the parent process! My child's PID is %d\n", pid);
}
printf("This message is printed by both the parent and the child.\n");
\end{lstlisting}

The convention for the return value of \texttt{fork} is:
\begin{itemize}
  \item negative $\Rightarrow$ failure;
  \item zero $\Rightarrow$ ``you're the child;''
  \item positive $\Rightarrow$ ``you're the parent, and this is your
    child's PID.''
\end{itemize}

\textbf{Caveat for \texttt{pid2.c}:} the parent doesn't
\texttt{wait()} for the child here, so the child can outlive the
parent (briefly) and become a ``zombie'' until reaped by init.
Always pair your forks with a \texttt{wait()} in real code.

\texttt{pid3.c} adds the \texttt{wait}:

\begin{lstlisting}[language=C]
else if(pid == 0){
  printf("Hey, give me a number, will ya?\n");
  int num;
  int success = scanf("%d",&num);
  return success ? 0 : 1;
}
else {
  wait(&childReturn);
  printf("Here's the int returned by childReturn: %d", childReturn);
}
\end{lstlisting}

But there's a gotcha: \texttt{wait}'s status output is not just the
exit code --- it's a packed integer where the exit code is in the
upper byte and the lower bits encode signals etc. So the printed
value isn't the child's return value; it's that value shifted left 8
bits.

\texttt{pid4.c} unpacks it manually:

\begin{lstlisting}[language=C]
wait(&childReturn);
// WEXITSTATUS(int) => (int >> 8) & 255
printf("Here's the int returned by childReturn: %d\n",
       (childReturn >> 8) & 255);
\end{lstlisting}

\texttt{noparent.c} uses the macro instead of the bit-fiddle:

\begin{lstlisting}[language=C]
wait(&childReturn);
if(WEXITSTATUS(childReturn) == 0){
  printf("Thanks for being nice to my child!\n");
}
else{
  printf("They've massacred my boy!\n");
}
\end{lstlisting}

Always use \texttt{WEXITSTATUS} (and friends, like
\texttt{WIFEXITED}, \texttt{WTERMSIG}) rather than rolling your own
shifts. They handle every corner case the kernel cares about.

\subsection*{Signals}

A signal is an asynchronous interrupt delivered to a process. The
kernel maintains a per-process table of handlers; when a signal
arrives, the kernel calls the corresponding handler. Useful signals:
\texttt{SIGINT} (Ctrl-C), \texttt{SIGTERM} (please die),
\texttt{SIGKILL} (die now, can't be caught), \texttt{SIGUSR1} /
\texttt{SIGUSR2} (your own use), \texttt{SIGSEGV},
\texttt{SIGALRM}.

\texttt{sig1.c} catches Ctrl-C:

\begin{lstlisting}[language=C]
void handler(int sig){
  (void)sig;
  printf("Caught that sigint!\n");
  exit(0);
}

int main(){
  printf("We have a ctrl-c handler here! I.E. SIGINT (%d)\n",SIGINT);
  signal(SIGINT, handler);
  while(1){
    printf("Beep boop\n");
    sleep(1);
  }
}
\end{lstlisting}

\texttt{sig2.c} makes the handler run multiple times before
exiting, with a countdown:

\begin{lstlisting}[language=C]
volatile sig_atomic_t num = 3;

void handler(int sig){
  (void)sig;
  if(num > 0){
    printf("Caught that sigint!\n");
    printf("Explosion in...%d\n",num);
    num--;
  }
  else{
    printf("Boom!\n");
    exit(0);
  }
}
\end{lstlisting}

The \texttt{volatile sig\_atomic\_t} matters: the variable is
modified asynchronously by the handler, so the compiler can't cache
its value in a register, and we need a type that's guaranteed to be
updated atomically with respect to signal delivery.

\texttt{tag.c} is the one I had the most fun with --- a parent and
child take turns sending each other signals:

\begin{lstlisting}[language=C]
volatile sig_atomic_t num = 0;

void handler1(int sig){ /* child */
  (void)sig;
  printf("I've been tagged by my parent %d times\n",num);
  num++;
}
void handler2(int sig){ /* parent */
  (void)sig;
  printf("I've been tagged by my child %d times\n",num);
  num++;
}

int main(){
  signal(SIGUSR1,handler1);
  signal(SIGUSR2,handler2);
  int pid = fork();
  if(pid < 0) return 1;
  else if(pid == 0){
    while(num < 5){
      kill(getppid(),SIGUSR2);
      sleep(1);
    }
  }
  else{
    while(num < 5){
      kill(pid,SIGUSR1);
      sleep(1);
    }
    kill(pid,SIGTERM);
    wait(0);
  }
}
\end{lstlisting}

The point here: signal delivery is an inter-process communication
primitive, not just a way for the kernel to tell you about errors.
Note the \texttt{wait(0)} at the end --- the parent must reap the
child after killing it, otherwise we leak a zombie.

\section{Lecture 8 --- Signals Continued, Threads Up Next}

A short, dense day on signals. We picked up the parent/child fight
that I'd promised and looked at signal masking, then started
threading at the very end.

\subsection*{Alarms}

\texttt{alarms.c} demonstrates \texttt{SIGALRM}, which the kernel
sends after a delay you request:

\begin{lstlisting}[language=C]
void alarm_handler(int sig){
  (void)sig;
  printf("The bells have been rung!\n");
}

int main(){
  signal(SIGALRM, alarm_handler);
  alarm(5);
  printf("We slumber\n");
  pause();        // sleep until any signal arrives
  printf("We have awoken!\n");
  return 0;
}
\end{lstlisting}

\texttt{pause()} is great for this: it puts the process to sleep
until \emph{any} signal arrives, then wakes it up.

\subsection*{The \texttt{rand()} fun fact}

\texttt{randtest.c} just calls \texttt{rand()} ten times in a row:

\begin{lstlisting}[language=C]
for(int i=0;i<10;i++){
  printf("Here's a random number %d\n",rand());
}
\end{lstlisting}

If you run it twice you'll notice it produces the same numbers each
time. \texttt{rand()} is a deterministic pseudorandom sequence,
seeded from a default value. To get different output you need
\texttt{srand(time(NULL))} (or something more entropic) at startup.

\subsection*{Signal masking and unmaskable signals}

\texttt{stopTest.c} tries to block \texttt{SIGSTOP}:

\begin{lstlisting}[language=C]
sigset_t blocks, prevs;
sigemptyset(&blocks);
sigaddset(&blocks,SIGSTOP);
sigprocmask(SIG_BLOCK,&blocks,&prevs);
while(1){
  printf("Here I am!: %d\n",getpid());
  sleep(1);
}
\end{lstlisting}

The point: \texttt{SIGSTOP} (the one Ctrl-Z sends, kind of, but
actually that's \texttt{SIGTSTP}) and \texttt{SIGKILL} can never be
blocked, ignored, or caught. The kernel forces them through. Try it
and you'll find Ctrl-Z still suspends the program.

\texttt{stopTest2.c} makes the same point with \texttt{SIGTSTP}, the
``terminal stop'' signal that Ctrl-Z actually sends:

\begin{lstlisting}[language=C]
void handler(int sig){
  (void)sig;
  printf("Mwhahaha\n");
}
int main(){
  signal(SIGTSTP,handler);
  while(1){
    printf("Here I am!: %d\n",getpid());
    sleep(1);
  }
}
\end{lstlisting}

This one \emph{can} be caught --- press Ctrl-Z and you get
``Mwhahaha'' instead of suspension. To kill the program you'll need
Ctrl-C (\texttt{SIGINT}) or \texttt{kill -9} (\texttt{SIGKILL}).

\subsection*{Parent/child signal fight}

\texttt{signalFight.c} is the whole-page program: parent and child
take turns sending each other \texttt{SIGUSR1}, randomly damaging
each other, until one runs out of HP.

\begin{lstlisting}[language=C]
volatile sig_atomic_t hp = 50;
volatile sig_atomic_t stillFighting = 1;
volatile sig_atomic_t won = 1;

void hitHandler(int sig){
  (void)sig;
  if(hp > 0){
    int d = rand()%5+1;
    char str[50];
    sprintf(str,"I, # %d, have been hit! I took %d damage!\n",getpid(),d);
    write(STDOUT,str,strlen(str));
    hp = hp - d;
    if(hp <= 0){ stillFighting = 0; won = 0; }
  }
}

void fightLoop(int enemy){
  while(stillFighting){
    char str[50];
    sprintf(str,"I, mr. %d, have %d hp left\n",getpid(),hp);
    write(STDOUT,str,strlen(str));
    kill(enemy,SIGUSR1);
    int rest;
    do { rest = sleep(rand()%4+1); } while(rest > 0);
  }
  kill(enemy,SIGUSR2);
}
\end{lstlisting}

Two crucial details inside the handler:
\begin{itemize}
  \item It uses \texttt{write()} on \texttt{STDOUT} rather than
    \texttt{printf}, because \texttt{printf} is not async-signal-safe
    --- if a signal interrupts the middle of a \texttt{printf} and
    your handler also calls \texttt{printf}, you can deadlock or
    corrupt output. \texttt{write} is safe.
  \item It does the formatting with \texttt{sprintf} into a local
    buffer first. (\texttt{sprintf} isn't strictly safe either, but
    it's much less risky than \texttt{printf} since it doesn't touch
    the stdio lock.)
  \item The \texttt{do/while(rest > 0)} loop around \texttt{sleep}
    handles the case where \texttt{sleep} is interrupted by a
    signal --- it returns the unslept time, and we keep going until
    we've actually slept the full duration.
\end{itemize}

The seeding ``\texttt{srand(time(0) \^{} getpid())}'' is a small
trick to keep the parent and child from getting identical random
sequences.

There's a commented-out alternative using \texttt{sigaction} with
\texttt{SA\_RESTART} at the top of \texttt{main} --- worth
internalizing as the more robust modern API. \texttt{signal()} is
fine for quick demos but \texttt{sigaction} gives you control over
exactly which other signals are masked while your handler runs, and
lets you get automatic syscall restart.

We'll get into pthreads next time.

\section{Lecture 9 --- Threads}

Threads are like processes that share memory. You get most of the
parallelism benefit and lose most of the isolation cost --- which is
also where the danger comes from. Today: how to start them, how to
wait for them, how to pass arguments and get results, and the first
mutex.

For more depth, \texttt{concurrencyGuide.org} in the repo is the
canonical reference for everything we're doing this week and next.

\subsection*{The basic shape}

A thread is started with \texttt{pthread\_create}, which takes a
thread handle, attribute pointer (usually \texttt{NULL}), function
pointer, and an opaque \texttt{void*} argument:

\begin{lstlisting}[language=C]
int pthread_create(pthread_t* thread, pthread_attr_t* attr,
                   void* (*func)(void*), void* arg);
\end{lstlisting}

The thread function takes a \texttt{void*} and returns a
\texttt{void*}. That \texttt{void*} is C's polymorphism: the caller
casts whatever they want into it on the way in, and the thread casts
it back to its real type inside.

\texttt{threads1.c}:

\begin{lstlisting}[language=C]
void* func1(void* arg){
  (void)arg;
  printf("Hi I'm a thread!\n");
  return NULL;
}

int main(){
  pthread_t thread1, thread2;
  pthread_create(&thread1,NULL,func1,NULL);
  pthread_create(&thread2,NULL,func1,NULL);
  pthread_join(thread1,NULL);
  pthread_join(thread2,NULL);
  return 0;
}
\end{lstlisting}

\texttt{pthread\_join} blocks until the named thread finishes ---
roughly the threading equivalent of \texttt{wait}. The second
argument is a place to stash the return value, or \texttt{NULL} if
you don't want it.

If you forget to join your threads, \texttt{main} can return out
from under them and the whole process gets cleaned up; sometimes
this looks like the threads ``didn't run.'' Always join.

\subsection*{Reading return values}

\texttt{threads2.c} captures the return value:

\begin{lstlisting}[language=C]
int* ret1; int* ret2;
pthread_create(&thread1,NULL,func1,NULL);
pthread_create(&thread2,NULL,func1,NULL);
pthread_join(thread1,(void**)&ret1);
pthread_join(thread2,(void**)&ret2);
printf("The value of ret1 is: %p\n",(void*)ret1);
\end{lstlisting}

The double-pointer is because we want \texttt{pthread\_join} to
write \emph{into} our \texttt{ret1}: it needs the address of our
pointer.

\texttt{threads3.c} actually returns something nontrivial:

\begin{lstlisting}[language=C]
void* func1(void* arg){
  (void)arg;
  int* sleepPointer = malloc(sizeof(int));
  *sleepPointer = rand()%5+1;
  sleep(*sleepPointer);
  return sleepPointer;
}
\end{lstlisting}

\textbf{Important point:} you cannot return a pointer to a stack
local from a thread function (or any function). The stack frame is
gone the moment the function returns. \texttt{malloc} the value if
you need it to outlive the thread; the joiner is responsible for
\texttt{free}-ing it.

\subsection*{Passing arguments in}

\texttt{threads4.c} uses the argument slot:

\begin{lstlisting}[language=C]
void* func1(void* arg){
  int sleepAmount = *(int*) arg;
  sleep(sleepAmount);
  printf("[Yawn] I slept for %d seconds\n",sleepAmount);
  return NULL;
}

int main(){
  int sleepArg1 = rand()%6 + 1;
  int sleepArg2 = rand()%7 + 2;
  pthread_create(&thread1,NULL,func1,&sleepArg1);
  pthread_create(&thread2,NULL,func1,&sleepArg2);
  ...
}
\end{lstlisting}

The threads dereference the pointer they were given. Note that the
\texttt{sleepArg1} and \texttt{sleepArg2} variables are local to
\texttt{main} but \texttt{main} doesn't return until after the joins,
so they live long enough.

\subsection*{Mutexes}

When two threads touch the same variable, you have a race. The
canonical illustration is incrementing a shared counter:

\texttt{threadmutex.c}:

\begin{lstlisting}[language=C]
int ourCounter = 0;
pthread_mutex_t mutex;

void* threadCounter(void* arg){
  (void)arg;
  pthread_mutex_lock(&mutex);
  int temp = ourCounter;
  sleep(rand()%3+1);
  ourCounter = temp+1;
  pthread_mutex_unlock(&mutex);
  return NULL;
}

int main(){
  pthread_t threads[10];
  pthread_mutex_init(&mutex,NULL);
  for(int i=0; i < 10; i++){
    pthread_create(&threads[i],NULL,threadCounter,NULL);
  }
  for(int i=0; i < 10; i++){
    pthread_join(threads[i],NULL);
  }
  pthread_mutex_destroy(&mutex);
  printf("The value of ourCounter is: %d\n",ourCounter);
}
\end{lstlisting}

The \texttt{int temp = ourCounter; sleep(\ldots); ourCounter = temp +
1} pattern is a deliberately bad idea --- it widens the race window
so much that without the mutex you'd get garbage almost every time.
With the mutex, only one thread can be inside the
\texttt{lock}/\texttt{unlock} block at a time, so the increments
serialize and you get exactly 10.

The mutex lifecycle: \texttt{pthread\_mutex\_init} once at startup,
\texttt{lock}/\texttt{unlock} pairs around critical sections,
\texttt{pthread\_mutex\_destroy} once at the end.

\section{Lecture 10 --- Mutexes Reviewed and Condition Variables}

Today we cleaned up the mutex example, then introduced condition
variables, then ran a producer/consumer demo to see them in action.

\subsection*{Mutex review: bundling lock with the data it protects}

\texttt{mutexreview.c} fixes a smell from last time: instead of one
global mutex protecting one global counter, we bundle the mutex
together with the counter into a struct.

\begin{lstlisting}[language=C]
struct protectedCounter{
  int counter;
  pthread_mutex_t mutex;
};

void* threadCounter(void* arg){
  struct protectedCounter* c = (struct protectedCounter*) arg;
  pthread_mutex_lock(&c->mutex);
  sleep(rand()%3+1);
  c->counter = c->counter+1;
  pthread_mutex_unlock(&c->mutex);
  return NULL;
}
\end{lstlisting}

Now you can have multiple independent counters, each with its own
mutex, and the threading code doesn't need to know which one it's
working on:

\begin{lstlisting}[language=C]
struct protectedCounter c1, c2;
c1.counter = 0; pthread_mutex_init(&c1.mutex,NULL);
c2.counter = 0; pthread_mutex_init(&c2.mutex,NULL);

for(int i=0; i < 10; i++){
  if(i % 2 == 1){
    pthread_create(&threads[i],NULL,threadCounter,&c1);
  } else {
    pthread_create(&threads[i],NULL,threadCounter,&c2);
  }
}
\end{lstlisting}

This is the design pattern: \emph{the mutex lives next to the data
it protects, and they're passed together everywhere.}

\subsection*{Condition variables}

A mutex protects shared state from concurrent modification. A
\emph{condition variable} lets a thread wait for shared state to
reach some condition. The two are designed to work together:

\begin{itemize}
  \item \texttt{pthread\_cond\_init} / \texttt{pthread\_cond\_destroy}
    --- lifecycle.
  \item \texttt{pthread\_cond\_wait(\&cond, \&mutex)} --- atomically
    unlock the mutex, sleep until signaled, re-acquire the mutex
    before returning.
  \item \texttt{pthread\_cond\_signal(\&cond)} --- wake one waiter.
  \item \texttt{pthread\_cond\_broadcast(\&cond)} --- wake all
    waiters.
  \item \texttt{pthread\_cond\_timedwait} --- like wait, but with a
    timeout.
\end{itemize}

The trick that makes condition variables work is the \emph{atomic}
unlock-and-sleep in \texttt{wait}. Without it, you'd have a window
between unlocking the mutex and going to sleep where a signal could
arrive and be missed.

\subsection*{The bespoke producer/consumer}

\texttt{condbespoke1.c} has two threads: a reader that asks for two
numbers and an adder that sums them, signaling each other when input
is ready.

\begin{lstlisting}[language=C]
pthread_mutex_t mut;
pthread_cond_t enoughInputs;
int numInputs = 0;
int inputArr[2];

void* adder(void* arg){
  pthread_mutex_lock(&mut);
  while(numInputs < 2){
    pthread_cond_wait(&enoughInputs,&mut);
  }
  int* num = (int*)arg;
  *num = inputArr[1] + inputArr[0];
  pthread_mutex_unlock(&mut);
  return NULL;
}

void* reader(void* arg){
  (void)arg;
  pthread_mutex_lock(&mut);
  while(numInputs < 2){
    printf("Enter a number:\n");
    scanf("%d",&inputArr[numInputs]);
    numInputs++;
  }
  pthread_cond_signal(&enoughInputs);
  pthread_mutex_unlock(&mut);
  return NULL;
}
\end{lstlisting}

Two patterns to commit to memory:
\begin{itemize}
  \item \emph{Always wait in a loop:} \texttt{while(numInputs < 2)}
    around the \texttt{pthread\_cond\_wait}, not \texttt{if}.
    Spurious wakeups are allowed by the standard, and you might
    re-check after \texttt{wait} returns and find the condition
    isn't actually satisfied yet.
  \item \emph{Always hold the mutex when checking the condition or
    calling wait/signal.} The mutex is what protects the variable
    that the condition is testing.
\end{itemize}

\subsection*{The classic producer/consumer}

\texttt{condtest.c} is the textbook example with a bounded
ring-buffer:

\begin{lstlisting}[language=C]
#define BUFFER_SIZE 5
int buffer[BUFFER_SIZE];
int count = 0, in = 0, out = 0;

pthread_mutex_t mutex;
pthread_cond_t cond_not_empty;
pthread_cond_t cond_not_full;

void* producer(void* arg){
  ...
  pthread_mutex_lock(&mutex);
  while (count == BUFFER_SIZE) {
    pthread_cond_wait(&cond_not_full, &mutex);
  }
  buffer[in] = item;
  in = (in + 1) % BUFFER_SIZE;
  count++;
  pthread_cond_signal(&cond_not_empty);
  pthread_mutex_unlock(&mutex);
  ...
}

void* consumer(void* arg){
  ...
  pthread_mutex_lock(&mutex);
  while (count == 0) {
    pthread_cond_wait(&cond_not_empty, &mutex);
  }
  int item = buffer[out];
  out = (out + 1) % BUFFER_SIZE;
  count--;
  pthread_cond_signal(&cond_not_full);
  pthread_mutex_unlock(&mutex);
  ...
}
\end{lstlisting}

The producer blocks if the buffer is full, the consumer blocks if
it's empty, and each one signals the other after every operation.
Two condition variables, one mutex --- the mutex protects all the
shared state (\texttt{buffer}, \texttt{count}, \texttt{in},
\texttt{out}) and the conditions describe ``the state has reached
something interesting.''

We talked briefly about semaphores too. A semaphore is like a
generalized mutex: it has a counter, and \texttt{sem\_wait}
decrements it (blocking if it's zero), \texttt{sem\_post}
increments it. With initial count 1, a semaphore behaves like a
mutex; with higher counts you can let $N$ threads through. Unlike
mutexes, semaphores can be shared across processes when allocated
in shared memory, which we'll see again at the end of the term.

\section{Lecture 11 --- File I/O in C}

We're done with concurrency for now (mostly). Today: how to read
and write files, both at the syscall level and through the
\texttt{stdio.h} \texttt{FILE*} layer.

\subsection*{Reading a file with raw syscalls}

\texttt{openFile1.c} uses \texttt{open}/\texttt{read}/\texttt{write}
from \texttt{<unistd.h>} --- the same syscalls we used in assembly:

\begin{lstlisting}[language=C]
int echoFile(char* fName, bool showSize){
  int fd = open(fName, O_RDONLY);
  char buffer[1024] = {0};
  ssize_t bytesRead = 0;

  if(showSize){
    struct stat fileStat;
    lstat(fName, &fileStat);
    printf("The file %s is %ld bytes long...\n", fName,
           (long)fileStat.st_size);
  }

  bytesRead = read(fd, buffer, sizeof(buffer));
  while(bytesRead > 0){
    write(STDOUT_FILENO, buffer, bytesRead);
    bytesRead = read(fd, buffer, sizeof(buffer));
  }
  close(fd);
  return 0;
}
\end{lstlisting}

The shape is: \texttt{open} returns a file descriptor (an int),
\texttt{read} fills your buffer, \texttt{write} drains it, and
you loop until \texttt{read} returns 0 (EOF) or a negative value
(error). Then \texttt{close} the descriptor.

\texttt{lstat} (or \texttt{stat}) gives you metadata --- size,
permissions, modification time, etc. \texttt{lstat} differs from
\texttt{stat} in how it handles symlinks: \texttt{lstat} returns
info about the link itself, \texttt{stat} follows the link.

The driver dispatches on a \texttt{-s} flag:

\begin{lstlisting}[language=C]
if(strcmp(argv[1],"-s") == 0){
  echoFile(argv[2],true);
}
else{
  echoFile(argv[1],false);
}
\end{lstlisting}

\subsection*{The \texttt{FILE*} interface}

\texttt{openFile2.c} uses the higher-level \texttt{stdio.h} API ---
\texttt{fopen} returns a \texttt{FILE*}, and you read line-by-line
with \texttt{fgets}:

\begin{lstlisting}[language=C]
FILE* ourFile = fopen(argv[1],"r");
char** lines = malloc(fileSize * sizeof(char*));
int numLines = 0;
lines[0] = malloc(lineSize * sizeof(char));
while(fgets(lines[numLines],lineSize,ourFile)){
  numLines++;
  lines[numLines] = malloc(lineSize * sizeof(char));
}
free(lines[numLines]);
\end{lstlisting}

The pattern is ``allocate a fresh buffer, read into it, advance the
counter, allocate the next one'' --- you end up with one extra
allocation when \texttt{fgets} returns NULL, which is why we
\texttt{free(lines[numLines])} at the end.

\subsection*{Building a tiny line editor}

\texttt{lineEditor.c} is the bigger payoff for the day --- a
\texttt{ed}-style text editor that loads a file, lets you insert /
delete / edit lines, and writes the result back when you quit.

\begin{lstlisting}[language=C]
void cleanUp(FILE* f, char** ls, int nL){
  ftruncate(fileno(f),0);
  fseek(f, 0, SEEK_SET);
  for(int i=0; i < nL; i++){
    fputs(ls[i],f);
    free(ls[i]);
  }
  free(ls);
  fclose(f);
}
\end{lstlisting}

The \texttt{ftruncate(fileno(f),0)} + \texttt{fseek(f,0,SEEK\_SET)}
combo is how you wipe a file and rewrite it from scratch:
\texttt{ftruncate} resizes the file to zero bytes, \texttt{fseek}
puts the cursor back at the start. Then we just \texttt{fputs} each
line back out.

\begin{lstlisting}[language=C]
void insLine(int line, char** ls, int* nL){
  char* newLine = malloc(linesize*sizeof(char));
  char c;
  printf("New text to insert at line %d:\n",line);
  while ((c = getchar()) != '\n' && c != EOF); // clear input buffer
  fgets(newLine, linesize, stdin);
  for(int i = *nL; i > line; i--){
    ls[i] = ls[i-1];
  }
  ls[line] = newLine;
  (*nL)++;
}
\end{lstlisting}

The \texttt{while ((c = getchar()) != '\textbackslash n' \&\& c !=
EOF);} pattern is a recurring annoyance in C: after a
\texttt{scanf("\%d",\&x)} reads a number, the trailing newline is
still in the input buffer, and the next \texttt{fgets} would slurp
just that newline and return immediately. So we eat the rest of
the line first.

The main loop is a switch over a menu choice; nothing fancy, but
it's a reminder that ``a real program'' is mostly bookkeeping
around the small interesting bits. \texttt{foop.txt} in the
directory is a sample file to experiment on.

\section{Lecture 12 --- Sockets, an Introduction}

Today (no agenda file in the directory --- I was running short on
time) we wrote our way through the socket API for TCP, ending with
a tiny one-page web server. The repo's \texttt{socketsGuide.org}
covers the whole API in much more depth if you want to keep going.

\subsection*{The basic server pattern}

The TCP server lifecycle is:
\begin{enumerate}
  \item \texttt{socket(AF\_INET, SOCK\_STREAM, 0)} --- get a socket
    file descriptor.
  \item Fill in a \texttt{struct sockaddr\_in} with family, address,
    port (port and address need to be in \emph{network byte order}
    --- use \texttt{htons} for the port, \texttt{INADDR\_ANY} for
    ``any address'').
  \item \texttt{bind} the socket to that address.
  \item \texttt{listen} on it (with a backlog size).
  \item \texttt{accept} returns a fresh socket fd for each incoming
    connection.
  \item \texttt{read}/\texttt{write} on that fd.
  \item \texttt{close} the per-connection fd when done.
\end{enumerate}

\subsection*{\texttt{sock1.c}: send a single greeting}

\begin{lstlisting}[language=C]
int port = 8080;
int server_fd, new_socket;
struct sockaddr_in address;
int addrlen = sizeof(address);
char buffer[1024] = "Hey there y'all\n";

server_fd = socket(AF_INET, SOCK_STREAM,0);
address.sin_family = AF_INET;
address.sin_addr.s_addr = INADDR_ANY;
address.sin_port = htons(port);

bind(server_fd, (struct sockaddr*)&address, sizeof(address));
listen(server_fd, 3);

new_socket = accept(server_fd, (struct sockaddr*)&address,
                    (socklen_t*)&addrlen);
write(new_socket, buffer, strlen(buffer));

close(new_socket);
close(server_fd);
\end{lstlisting}

To test it, run the program in one terminal, then in another do
\texttt{nc localhost 8080} (netcat). You'll see ``Hey there
y'all''.

\subsection*{\texttt{sock2.c}: an echo server for one client}

\texttt{sock2.c} adds a read/write loop so the server bounces
whatever it gets back to the client:

\begin{lstlisting}[language=C]
new_socket = accept(server_fd, (struct sockaddr*)&address,
                    (socklen_t*)&addrlen);
while(1) {
  memset(buffer, 0, 1024);
  int valread = read(new_socket, buffer, 1024);
  if (valread > 0) {
    printf("Received: %s", buffer);
    write(new_socket, buffer, strlen(buffer));
  } else {
    break;
  }
}
\end{lstlisting}

This handles a single client; when they hang up, \texttt{read}
returns 0 and we break out. The server then exits. Not super useful,
but a step toward the real thing.

\subsection*{\texttt{sock3.c}: handle multiple clients sequentially}

\texttt{sock3.c} wraps the \texttt{accept}/echo loop in an outer
loop bounded by \texttt{maxClients = 5}:

\begin{lstlisting}[language=C]
while(clients < maxClients){
  new_socket = accept(server_fd, (struct sockaddr*)&address,
                      (socklen_t*)&addrlen);
  clients++;
  while(1) {
    memset(buffer, 0, 1024);
    int valread = read(new_socket, buffer, 1024);
    if (valread > 0) {
      printf("Received: %s", buffer);
      write(new_socket, buffer, strlen(buffer));
    } else {
      break;
    }
  }
}
\end{lstlisting}

This is sequential: the server services one client to completion
before accepting the next. To handle multiple clients
concurrently, you'd \texttt{fork} after \texttt{accept} (one process
per client), or spawn a thread per client, or use \texttt{select}
/ \texttt{epoll} to multiplex on a single thread. We'll see those
patterns in the next lecture and in the guide.

\subsection*{\texttt{websock.c}: a minimal HTTP server}

The grand finale: a server that speaks HTTP/1.1 well enough to
return ``Hello, World'' to a browser:

\begin{lstlisting}[language=C]
while (1) {
  memset(buffer,0,BUFFER_SIZE);
  new_socket = accept(server_fd, (struct sockaddr*)&address,
                      (socklen_t*)&addrlen);
  read(new_socket, buffer, BUFFER_SIZE);
  printf("Received request:\n%s\n", buffer);

  char *response = "HTTP/1.1 200 OK\r\n"
                   "Content-Type: text/html\r\n"
                   "Content-Length: 46\r\n"
                   "\r\n"
                   "<html><body><h1>Hello, World!</h1></body></html>";

  write(new_socket, response, strlen(response));
  close(new_socket);
}
\end{lstlisting}

HTTP is just lines of text over TCP. The status line, then headers,
then a blank line (\texttt{\textbackslash r\textbackslash n\textbackslash r\textbackslash n}),
then the body. The \texttt{Content-Length: 46} header is the byte
count of the body --- if you change the body, you have to update
that number, otherwise the browser will hang waiting for bytes that
aren't coming. Run this and visit \texttt{http://localhost:8080} in
a browser; ignore the server log of the request that flies by.

\section{Lecture 13 --- IPC: Pipes, Named Pipes, and Shared Memory}

The last regular lecture: a tour of the various ways processes can
communicate without going through sockets. None of this is on an
assignment, it's just fun.

\subsection*{Anonymous pipes}

\texttt{pipe1.c}: a parent and child share a pipe, child writes,
parent reads.

\begin{lstlisting}[language=C]
int pipefd[2];   // [0] read end, [1] write end
char write_msg[] = "Hello from the child process!";
char read_msg[100];

if(pipe(pipefd) == -1) { perror("Pipe failed"); return 1; }

pid_t pid = fork();

if (pid == 0) {              // child
  close(pipefd[0]);          // child doesn't read
  write(pipefd[1], write_msg, strlen(write_msg) + 1);
  close(pipefd[1]);
}
else {                       // parent
  close(pipefd[1]);          // parent doesn't write
  read(pipefd[0], read_msg, sizeof(read_msg));
  printf("Parent received: %s\n", read_msg);
  close(pipefd[0]);
}
\end{lstlisting}

The pattern: \texttt{pipe()} fills in a 2-element array of fds,
then you \texttt{fork}, then each side closes the end it doesn't
need. Pipes are unidirectional --- if you need bidirectional
communication you make two pipes.

(\textbf{Caveat:} this code doesn't \texttt{wait()} for the
child. Add a \texttt{wait(NULL)} at the end of the parent branch in
real code.)

\subsection*{Two pipes, both ways: \texttt{tempconverter.c}}

\texttt{tempconverter.c} demonstrates bidirectional communication:
parent sends Fahrenheit values to child, child converts and sends
back Celsius:

\begin{lstlisting}[language=C]
int parent_to_child[2];  // Parent writes, child reads
int child_to_parent[2];  // Child writes, parent reads

pipe(parent_to_child);
pipe(child_to_parent);

pid_t pid = fork();
if (pid > 0) {  // parent
  close(parent_to_child[READ_END]);
  close(child_to_parent[WRITE_END]);
  for (int i = 0; i < count; i++) {
    sprintf(buffer, "%.1f", temperatures[i]);
    write(parent_to_child[WRITE_END], buffer, strlen(buffer)+1);
    read(child_to_parent[READ_END], buffer, BUFFER_SIZE);
    printf("Parent: Received result: %s\n", buffer);
  }
  strcpy(buffer, "EXIT");
  write(parent_to_child[WRITE_END], buffer, strlen(buffer)+1);
  wait(NULL);
}
\end{lstlisting}

Two pipes, each side closes the ends it doesn't use, and the
parent sends a sentinel \texttt{"EXIT"} message to tell the child
to clean up.

\subsection*{Named pipes (FIFOs)}

A \emph{named pipe} (a FIFO) is like a pipe with a path on the
filesystem. Two unrelated processes can both \texttt{open} the path
and talk to each other.

\texttt{namedPipesServer1.c} creates two FIFOs and reads from one /
writes to the other:

\begin{lstlisting}[language=C]
mkfifo("clientToServer", 0666);
mkfifo("serverToClient", 0666);

fdWrite = open("serverToClient", O_WRONLY);
fdRead  = open("clientToServer", O_RDONLY);

while(!stop_requested){
  memset(buffer, 0, bSize);
  int bytesRead = read(fdRead, buffer, bSize);
  if(bytesRead > 0){
    printf("Client: %s\n", buffer);
    write(fdWrite, buffer, bytesRead);
  }
  ...
}
\end{lstlisting}

\texttt{mkfifo} creates the file in the filesystem with the right
type. Note the open order matters: \texttt{open(O\_WRONLY)} blocks
until someone opens the other end for reading. The whole client
file is in \texttt{namedPipesClient1.c} --- it spawns a reader
thread to print whatever the server says, and reads stdin from the
main thread.

The chat server/client (\texttt{chat\_server.c} /
\texttt{chat\_client.c}) is a polished version of the same idea
with proper signal-handling for cleanup.

\subsection*{Shared memory + semaphores}

\texttt{shm\_open} creates a chunk of memory that can be mapped
into multiple processes' address spaces.

\texttt{sharedstruct.h} defines what they'll see:

\begin{lstlisting}[language=C]
struct sharedData {
  sem_t mutex;
  int payload1;
  int payload2;
};
#define sharedName "/shared_data_mem"
\end{lstlisting}

The semaphore lives \emph{inside} the shared region so both
processes can see it. \texttt{sharedMemprod.c}:

\begin{lstlisting}[language=C]
int sharedFd = shm_open(sharedName, O_CREAT | O_RDWR, 0666);
ftruncate(sharedFd, sizeof(struct sharedData));
struct sharedData* shared = mmap(NULL, sizeof(struct sharedData),
                                 PROT_READ | PROT_WRITE,
                                 MAP_SHARED, sharedFd, 0);
sem_init(&shared->mutex, 1, 1);   // pshared=1 -> cross-process
for(int i = 0; i < 10; i++){
  sem_wait(&shared->mutex);
  shared->payload1 = i;
  shared->payload2 = i*i;
  sem_post(&shared->mutex);
  sleep(1);
}
\end{lstlisting}

The \texttt{sem\_init(\&shared->mutex, 1, 1)}'s middle argument is
\emph{pshared}: 1 means ``shared between processes,'' 0 means
``shared between threads only.'' This is the bit that lets one
process wait on a semaphore another process posts to.

\texttt{sharedMemcons.c} is the consumer: it opens the same shared
region (without \texttt{O\_CREAT}), maps it, and reads from it
under the same semaphore. At the end it calls
\texttt{shm\_unlink(sharedName)} to remove the kernel-resident
object so a fresh run starts clean.

The bigger \texttt{shm\_creator.c} / \texttt{shm\_client.c} pair
implements a collaborative drawing board: the creator displays a
20x10 character grid, and any number of clients can connect and
randomly drop their character into the grid. Worth running with two
or three clients in different terminals.

\subsection*{Message queues}

\texttt{task\_distributor.c} and \texttt{worker.c} use SysV
message queues (\texttt{msgget}, \texttt{msgsnd}, \texttt{msgrcv}).
The distributor sends task messages tagged \texttt{TASK\_MSG = 1};
workers receive only that tag, do the (fake) work, and send
back-tagged \texttt{RESULT\_MSG = 2}. The queue acts as a typed,
multi-consumer mailbox in the kernel. We didn't dig deeply into
this one --- it's old API and most modern code reaches for sockets
or shared memory instead --- but it's a useful tool to know exists.

\section*{Appendix --- Struct Alignment and Optimization}

This was a half-lecture I slotted in for some of the loose ends I'd
been wanting to mention --- struct field alignment, common compiler
optimizations, and a quick aside about cache coherency. The files
are in \texttt{lecture-inf/}.

\subsection*{Why does this struct take more space than I expected?}

\texttt{structsize.c}:

\begin{lstlisting}[language=C]
struct Goofy1 { int num1; int num2; };          // 8 bytes
struct Goofy2 { int num1; double num2; };       // 16 bytes
struct Goofy3 { double num1; int num2; };       // 16 bytes
\end{lstlisting}

\texttt{Goofy1} is a clean 8 bytes, but \texttt{Goofy2} and
\texttt{Goofy3} are both 16. The reason is \emph{alignment}:
hardware reads memory in chunks aligned to the architecture's word
size, and to read a \texttt{double} the address needs to be a
multiple of 8.

In \texttt{Goofy2}, the \texttt{int} takes 4 bytes; the
\texttt{double} that follows it must start at offset 8, so the
compiler inserts 4 bytes of padding between them.

In \texttt{Goofy3}, the \texttt{double} starts at offset 0 (fine),
the \texttt{int} starts at offset 8 (4 bytes), and then the struct
itself must be padded out to 16 so that arrays of \texttt{Goofy3}
have each element properly aligned for its first \texttt{double}.

\texttt{structsize2.c} is a more elaborate example with multiple
arrays:

\begin{lstlisting}[language=C]
struct exA {  int int_field[3];  char char_field;  long array[2];  };
struct exB {  char char_field;   long array[2];    int int_field[3]; };
struct exC {  long array[2];     int int_field[3]; char char_field; };
\end{lstlisting}

The naive byte counts are all the same (29), but the actual
\texttt{sizeof}s differ because of where the padding lands. The
moral: \emph{if you care about size, sort fields from largest to
smallest}.

\subsection*{Loop unrolling}

For tight loops where the body is short, the cost of the comparison
and jump can be larger than the body itself. ``Unrolling'' means
doing several iterations' worth of work per loop pass:

\begin{lstlisting}[language=C]
// instead of:
for(int i = 0; i < 1000; i++){ arr[i] = i; }

// you might write (or the compiler might generate):
for(int i = 0; i < 1000; i += 5){
  arr[i]   = i;
  arr[i+1] = i+1;
  arr[i+2] = i+2;
  arr[i+3] = i+3;
  arr[i+4] = i+4;
}
\end{lstlisting}

This pairs nicely with the CPU's branch predictor and pipeline:
fewer branches, more in-flight arithmetic. Don't do this by hand
unless you have a real reason to --- the optimizer is better at
deciding when it pays off than you are.

\subsection*{Function inlining}

Calling a function costs a few instructions (push args, call,
return, restore). For functions whose body is shorter than the call
sequence, the compiler can substitute the body in place at the call
site. The \texttt{inline} keyword is a hint; \texttt{static inline}
and link-time optimization mostly take it from there.

\subsection*{Loop fusion}

Two loops over the same range can be combined into one:

\begin{lstlisting}[language=C]
// before:
for(int i=0; i < 1000; i++){ arr[i] = i*i; }
for(int i=0; i < 1000; i++){ printf("arr[%d] = %d", i, arr[i]); }

// after:
for(int i=0; i < 1000; i++){
  arr[i] = i*i;
  printf("arr[%d] = %d", i, arr[i]);
}
\end{lstlisting}

You save a loop's worth of overhead, and you also keep \texttt{arr[i]}
hot in cache between the write and the read.

\subsection*{Cache coherency}

The biggest performance win on modern hardware is being kind to the
cache. Consecutive accesses to nearby addresses are essentially free;
random accesses across megabytes are agonizingly slow. The classic
example is matrix iteration order:

\begin{lstlisting}[language=C]
int m[100][200];
int accum = 0;

// bad: column-major access of a row-major array
for(int j=0; j<200; j++){
  for(int i=0; i<100; i++){
    accum += m[i][j]*m[i][j];
  }
}

// good: row-major access of a row-major array
for(int i=0; i<100; i++){
  for(int j=0; j<200; j++){
    accum += m[i][j]*m[i][j];
  }
}
\end{lstlisting}

In C, \texttt{m[i][j]} and \texttt{m[i][j+1]} are adjacent in
memory; \texttt{m[i][j]} and \texttt{m[i+1][j]} are 200 ints apart.
The good version touches 200 contiguous ints before moving to the
next row, so each cache line gets fully used. The bad version
strides through memory and constantly evicts/reloads cache lines.
On a big matrix this is easily a 10x performance difference.

That's all for the term --- the assignments, posit tutorial, and
guides should round things out. See you in office hours.

\end{document}
